\section{Why the quadratic approximation stops at the fourth-order scale}
\label{p13:lower:section}

The quadratic reference retains the second factorial cumulant linearly.
Two independent pair contributions therefore leave a fourth-order term of
size $\alpha_N^2$. To turn this observation into a lower bound, one must
check that incompatible arithmetic configurations cannot cancel that term.
The half-exponent count supplies this missing sign information.
The formal term \(\alpha_N^2\Pi_{\mu_N}*\Delta^{*4}/8\)
is classical in Poisson--Charlier expansions
\cite{bkn2014,cekanaviciusvellaisamy2015}. The issue here is its
arithmetic size, the sign of the other fourth-order contributions,
and the passage from an unbounded polynomial test to the full norm.
Throughout this section $\mu=\mu_N$, $\alpha=\alpha_N$, and all marginals
and centered coefficients are those of the full admissible set.

\subsection{Cancellation of the disconnected pair terms}

Write
\[
 C_2^{\mathrm s}=\sum_{|e|=2}c(e),\qquad
 K_4=\sum_{\substack{|I|=4\\I\text{ connected}}}c(I),\qquad
 K_4^+=\sum_{\substack{|I|=4\\I\text{ connected},\ q(I)>0}}c(I).
\]
Compatibility makes every character on the relation space positive, so
$K_4^+\geq0$.

\begin{lemma}[The fourth factorial cumulant]
\label{p13:lower:quartic}
For every fixed critical band and every $\varepsilon>0$,
\begin{equation}
 \kappa_4=24K_4+O_{C,\varepsilon}(N^{-5/3+\varepsilon}),\qquad
 K_4=K_4^++O_{C,\varepsilon}(N^{-3/2+\varepsilon}).
 \label{p13:lower:quartic-identities}
\end{equation}
Consequently
\begin{equation}
 \kappa_4+3\alpha^2
 =3\alpha^2+24K_4^++O_{C,\varepsilon}(N^{-3/2+\varepsilon}).
 \label{p13:lower:positive-quartic}
\end{equation}
\end{lemma}

\begin{proof}
Expanding the first four factorial moments with the individual $p_i$
gives the exact identity
\begin{align}
 \kappa_4={}&24\sum_{|I|=4}c(I)-12(C_2^{\mathrm s})^2
 -24\sum_{|T|=3}\Bigl(\sum_{i\in T}p_i\Bigr)c(T)\notag\\
 &+24\sum_{e=\{i,j\}}(p_i^2+p_ip_j+p_j^2)c(e)
 -6\sum_i p_i^4.
 \label{p13:lower:exact-cumulant}
\end{align}
For example, first write
\[
 \E(1+t)^Z=\prod_i(1+p_it)
 \left[1+\sum_{2\leq|I|\leq4}
       \frac{c(I)t^{|I|}}{\prod_{i\in I}(1+p_it)}+O(t^5)\right],
\]
then take $24$ times the coefficient of $t^4$ in its formal logarithm.
This is a finite coefficient calculation; no analytic logarithm of the
true generating function is required.

Canonical direct decomposition gives
$\sum_{|I|=4}c(I)=K_4+D_{22}$, where $D_{22}$ sums
$c(e)c(f)$ over unordered disjoint pairs of edges whose spaces are direct.
In $(C_2^{\mathrm s})^2/2-D_{22}$ only repeated edges, edges sharing a
vertex, and disjoint edges with a nonzero intersection remain.
The positive propagation estimate, together with the product of the
marginals, implies
\[
 \sum_{y\ne x}|c(x,y)|\ll_{C,\eta}p_xN^\eta.
\]
Repeated edges and shared vertices consequently have total absolute
weight $O(p_*A_2N^\eta)=O(N^{-5/3+2\eta})$.

Apply the pair partition of Lemma~\ref{p13:activity:pair-partition} with
a small fixed parameter $\eta$. A product involving a weak or exceptional
edge has total weight at most
$A_2(A_{\mathrm w}+A_{\mathrm s})
\ll N^{-5/3+O(\eta)}$.
For two good edges from the remaining finite family of channels, fix the
first edge and the offset pattern of the second compression. There are
$NB^{O(1)}$ choices for the first edge. A nonzero intersection imposes
an exceptional target of fixed width, giving $N^{1/2+o(1)}$ choices for
the second compressed center by the elementary exceptional-target count
in Proposition~\ref{p1:arith:fixed-branching}. The proof of that count
uses only a fixed width and applies to this compression as well.
Each product has weight at most $N^{-10/3+o(1)}$. Hence this last cost is
$N^{-11/6+o(1)}$, and
\[
 D_{22}=(C_2^{\mathrm s})^2/2+O(N^{-5/3+\varepsilon}).
\]
The remaining terms of \eqref{p13:lower:exact-cumulant} are bounded by
$O(p_*A_3+p_*^2A_2+p_*^3\mu)$.
The triple estimate in Corollary~\ref{p13:third:activity} gives
$A_3\ll N^{-41/30+\eta}$, so each is smaller than the stated error.
This proves the first identity.

The incompatible part of Proposition~\ref{p13:transfer:fixed-activities}
at the fixed order \(r=4\) gives
\[
 |K_4-K_4^+|
 \le\sum_{\substack{|I|=4,\ I\text{ connected}\\q(I)=0}}|c(I)|
 \ll_{C,\varepsilon}N^{-3/2+\varepsilon}.
\]
Compatible coefficients are nonnegative. This proves the second identity
and then \eqref{p13:lower:positive-quartic}.
\end{proof}

\subsection{A Charlier polynomial detects the remaining mass}

\begin{theorem}[Optimal power for the exact quadratic reference]
\label{p13:lower:optimal}
For every $C>0$ and $\varepsilon>0$, there are
$c_{C,\varepsilon}>0$ and $N_0(C,\varepsilon)$ such that
\begin{equation}
 \|\Law(Z_N)-Q_{\mu_N,\alpha_N}\|_1
 \geq c_{C,\varepsilon}N^{-4/3-\varepsilon}
 \label{p13:lower:norm}
\end{equation}
whenever $N\geq N_0$ and $|L-\log_2N|\leq C$.
Together with Theorem~\ref{p13:transfer:reconstruction} and the triple
activity bound, this identifies the optimal power as $4/3$.
It does not assert an asymptotic constant or a logarithmic factor for
the remainder.
\end{theorem}

\begin{proof}
The arithmetic coefficient estimate gives
$\alpha\asymp_C LN^{-2/3}>0$ uniformly in the critical band.
The sign argument has the margin $3/2-4/3=1/6$.
Choose a fixed total error loss $\tau<1/6$ in
\eqref{p13:lower:positive-quartic}; then
\[
 \frac{N^{-3/2+\tau}}{\alpha^2}
 \ll_C L^{-2}N^{-1/6+\tau}\longrightarrow0.
\]
Positivity therefore yields
\[
 \kappa_4+3\alpha^2\geq (3-o_C(1))\alpha^2.
\]
The fourth Charlier polynomial is
\[
 H_4(z;\mu)=(z)_4-4\mu(z)_3+6\mu^2(z)_2-4\mu^3z+\mu^4
          =4![t^4]e^{-\mu t}(1+t)^z.
\]
Its integral against $Q_{\mu,\alpha}$ vanishes, because the factorial
generating function of that reference is
$e^{\mu t}(1+\alpha t^2/2)$. On the true law its expectation equals
$\kappa_4+3\alpha^2$. Thus, with
$R=\Law(Z_N)-Q_{\mu,\alpha}$,
\[
 \left|\int H_4\,dR\right|\gg_C L^2N^{-4/3}.
\]

The test is unbounded, so a moment argument is needed before it gives
a norm estimate. For every fixed integer $k>1$, the fixed moments from
Theorem~\ref{p1:fixed:moments} give
$\int |H_4(z;\mu)|^k\,d\Law(Z_N)(z)\ll_{C,k}1$.
The same bound holds for $|Q_{\mu,\alpha}|$, since $\mu$ stays in a
compact subset of $(0,\infty)$, $\alpha=O_C(1)$, and
\[
 |Q_{\mu,\alpha}|
 \leq\Pi_\mu+\frac{|\alpha|}{2}
       \Pi_\mu*(\delta_2+2\delta_1+\delta_0).
\]
Consequently H\"older's inequality for the positive measure $|R|$ gives
\[
 \left|\int H_4\,dR\right|
 \ll_{C,k}\|R\|_1^{1-1/k},\qquad
 \|R\|_1\gg_{C,k}(L^2N^{-4/3})^{k/(k-1)}.
\]
After $\varepsilon$ has been prescribed, choose the integer $k$ so that
$4/[3(k-1)]\leq\varepsilon$. It remains fixed as $N$ tends to infinity,
and proves \eqref{p13:lower:norm}. The upper bound follows from the
reconstruction and $A_3\ll N^{-41/30+\eta}$, since $41/30>4/3$.
\end{proof}

The obstruction is specific to the quadratic reference. The Hermite law
with the same first two factorial cumulants already contains the term
$\alpha^2\Pi_\mu*\Delta^{*4}/8$. The fourth Charlier test against that
law sees $\kappa_4$ instead of $\kappa_4+3\alpha^2$, so the lower bound
above does not establish an optimal Hermite approximation rate.
